Supóngase que $\Psi=\psi\otimes\varphi$, con $\psi,\varphi\in H_1$. Si Ana mide en una base $\{u_0,u_1\}$, el resultado $k$ tiene probabilidad $p_k=|\bk{u_k}{\psi}|^2$, y la onda de Bruno es
\[
  \Psi_k(s)=\bk{u_k}{\psi}\,\varphi(s).
\]
Tras normalizar, $\varphi_k$ es $\varphi/\|\varphi\|$ multiplicada por un factor de fase, para todo $k$ con $p_k\neq0$. En otras palabras, el estado de Bruno es siempre el mismo, sea cual sea la base y el resultado de Ana.

El paso del argumento de 2.4 que falla es la suposición de que el bit de Ana cambia el estado de Bruno. Con un estado producto se tiene
\[
  \rho_0=\rho_1=\varphi^{\otimes n},
\]
donde $\varphi^{\otimes n}:=\varphi(s_1)\cdots\varphi(s_n)$. Es decir, el mismo estado puro aparece en los dos casos. Entonces ningún procedimiento distingue los dos bits, y el error es $\frac12$ sin importar $n$.

Conviene notar que en el caso entrelazado de 2.3 la diferencia aparece solo al usar varias copias. Con una sola copia, el estado promedio de Bruno es $\frac12 I$ para los dos bits, pero $\rho_0$ y $\rho_1$ son mezclas distintas de estados puros, y al usar $n$ copias esa diferencia se hace visible. Con un estado producto, en cambio, $\rho_0$ y $\rho_1$ son el mismo estado puro desde el principio. Por tanto, lo que permite el argumento es el entrelazamiento de $\Psi$, que por 2.2 se caracteriza por $\det C\neq0$.
